\section{CPU TPC-H prepared-input certificates}

These are IO certificates, not compiler evidence or query annotations.  Let
$F$ be an immutable Parquet file, $R(F)$ its decoded rows, $H$ its heading,
$K\subseteq H$ a nonempty column set, and $P\subseteq H$ a loaded projection.
Preparation verifies keys against the exported file, not a declared TPC-H
schema constraint.  A full-heading certificate is the special case $K=H$.

\Rule{IO-CpuPrepareKey}
{K\ne\varnothing \quad K\subseteq H \quad
 |\pi_K R(F)|_{\mathrm{distinct}}=|R(F)| \quad
 \forall r\in R(F),k\in K.\ r[k]\ne\mathrm{NULL}}
{\operatorname{prepare}(F,K)=
 \operatorname{certificate}(\operatorname{sha256}(F),K)}

The preparer chooses the table's conventional primary-key columns when all
are present, otherwise the complete exported heading. One grouped traversal
checks duplicate keys and null key values together. Failure emits no key,
without trying another candidate or removing raw rows. Successful certificates
carry the explicit marker \code{set:"verified"}; unsuccessful records carry
\code{set:"uncertified"} and no verified keys. A manifest-only command may
write these records outside the immutable dataset.

\Rule{IO-CpuLoadKey}
{c=\operatorname{certificate}(h,K)\in M \quad c.\mathrm{set}=\mathrm{verified} \quad
 h=\operatorname{sha256}(F) \quad K\subseteq P \quad
 \operatorname{decode}|_K\ \text{is injective}}
{\operatorname{load}_P(F,M)=\pi_P R(F)\quad
 \text{without runtime duplicate elimination}}

Without these premises the existing duplicate-eliminating loader is used.
In particular, a key omitted by a projection and a full-heading certificate
for a wider heading do not authorize the shortcut. Missing, malformed,
stale or mismatched manifests do not establish a premise.  Certificates
are trusted preparation artifacts bound to bytes; a digest is not a
signature and an arbitrary user-authored assertion is not a verified key.
Host-built data has no certificate.  The decoder must preserve equality on
the verified key; lossy decoding cannot use this shortcut.
